\section{Conclusion}
\label{sec:conclusion}

Bounded addition and inversion can be implemented by a sparse resistive
network while retaining every source solution and every coordinate needed to
recover it. The resulting exact feasibility problem is $\ER$-complete even
under the simultaneous structural restrictions of
Theorem~\ref{thm:structural}. The model allows independently prescribed
voltage and injection bounds, including simultaneous singletons. Within
this model, the graph and coefficient restrictions do not suffice to make
exact feasibility easier than general existential-real decision. The
solution-preserving electrical realization is the main contribution;
the arithmetic and topological tools on which it builds are established.

The same correspondence explains the algebraic and numerical behavior of
feasible voltage sets. Rational equivalence realizes precisely the compact
basic closed sets over $\Q$, while semialgebraic homeomorphism realizes all
compact semialgebraic topologies. Unique voltages can have arbitrarily high
algebraic degree. Separately, the ordinary bounded-degree, fixed-data
construction gives infeasible networks with doubly exponentially small
injection residuals. This matches the general separation scale for that
class; no residual-preserving planarization is asserted. Exact decision,
tolerance-based numerical evidence, and certificates for a gap promise
consequently require different guarantees.

For AC models, the transfer depends on the meaning of an angle limit.
An explicit rational quadratic formula encodes the established zero-winding
condition using linearly many variables and predicates. Reactive balance
then gives both an exact equal-angle theorem and a quantitative stability
estimate. Principal-angle constraints alone allow winding states, whereas
the stated bus-angle boxes and size-dependent principal windows support
separate hardness transfers. These conclusions classify exact feasibility
for the specified models and delimit residual-based certification; they do
not predict typical solver performance or cover other operating models
that exclude the prescribed voltage and injection combinations used by the
reduction.
